\section{Protocol and reproducibility}
\label{app:protocol}

The frozen experiment families are:
\begin{itemize}[leftmargin=1.3em,itemsep=0.2em]
\item Primary: \nolinkurl{slowlab-v2-phase4-confirmatory-2026-09-13-a2}, with 928 complete and unique expected episode identities.
\item Noise extension: \nolinkurl{slowlab-v2-phase4-noise-robustness-2026-09-19-a2}, with 384 complete rows.
\item Post-primary model extension: \nolinkurl{slowlab-v2-phase5-model-extension-2026-09-20-a1}.
\end{itemize}
Each record stores requested and actual model identifiers, provider, generation seed, prompt
hashes, usage, cost, retry attempt, finish reason, format failures, infeasible submissions,
and the accepted event history.

The public numerical source for the current tables is
\nolinkurl{results/phase5_paper_summary_env2.0.0.json}. From an unpacked frozen data bundle,
\nolinkurl{scripts/reproduce_phase5_paper.py} rebuilds analyses, generated tables, figures,
claim checks, cross-reference checks, and the PDF. Private review notes are excluded from
the repository artifact.

\section{Reference policies}
Random spread maximises geometric coverage subject to the same facility validator as model
agents. Split-plot DOE balances whole-plot combinations across chambers and randomises
loop-level settings within chamber. Constraint-aware batch Bayesian optimisation fits an
ordinary GP to completed outcomes and proposes a feasible batch by sequential
fantasisation. Reader replay copies each accepted completed history to ordinary,
block-aware, and component-aware GPs without changing a design or observation.

\section{Deferred dynamic actions}
The current environment records time-stamped interim measurements but exposes no treatment
change, early termination, cleanup, residual-value, or asynchronous replacement action.
Adding one requires explicit transition dynamics and a matched oracle. Until then, claims
about the value of within-cycle observation are restricted to recommendation and next-round
behaviour.
